\documentclass[12pt]{article}
\usepackage[utf8]{inputenc}
\usepackage[T1]{fontenc}
\usepackage{hyperref}
\title{Tondi Songway Kiini Language Resources for K-12 Curriculum Development}
\author{Compiled by Claude Sonnet 5 (claude-sonnet-5) for Andr\'as Kornai}
\date{}
\begin{document}
\maketitle
\section{Introduction}
Tondi Songway Kiini (ISO 639-3 tst) is a Southern Songhay language of Mali, spoken in villages around Kikara, roughly 120 km west of Hombori. One source cites about 3,000 speakers (1998). It is unusual among Malian Songhay languages in having lexical tone. The search found no school materials in the language.
\section{Primary Educational Resources}
None were found in search results.
\section{Linguistic and Reference Materials}
Jeffrey Heath's 2005 reference grammar (CSLI Publications, about 450 pages) includes a Tondi Songway Kiini--English--French dictionary: \url{https://web.stanford.edu/group/cslipublications/cslipublications/site/157586505X.shtml}. Further overviews are at \url{https://en.wikipedia.org/wiki/Tondi_Songway_Kiini} and \url{https://websites.umich.edu/~jheath/Songhay%20languages.html}.
\section{Religious and Community Resources}
None were found in search results.
\section{Digital and Audio Resources}
The Endangered Languages Project has a catalogue entry: \url{https://endangeredlanguages.com/node/95674}. No audio or corpus was found.
\section{Curriculum Development Context}
Heath's grammar and dictionary are the only substantial resource found. A native-language curriculum would have to start from that description and from new work with the community. No primers, textbooks or school programmes were found.
\end{document}